\documentclass[aspectratio=169,9pt]{beamer}
\mode<presentation> 
\usepackage[utf8]{inputenc}
\usepackage{sansmathaccent} %Für equations
\usepackage{amsmath}
\usepackage{tikz}
\usepackage{multicol}
\usepackage{cancel}
\usepackage{relsize}
\usepackage{graphicx}
\usepackage[table]{xcolor}
\usepackage{fontawesome5}
\usepackage{enumitem}
\usepackage{setspace}
\usepackage[labelfont={bf},font={color=black!40,footnotesize},tablename={Tabelle},figurename={Abbildung}]{caption}
\usetikzlibrary{positioning, arrows.meta, calc}
\usefonttheme[onlymath]{serif}
%\makeatletter
%\makeatother

\usetikzlibrary{patterns, arrows.meta, calc}
\usepackage{pgfplots}
\pgfplotsset{compat=1.18}
\usetikzlibrary{intersections, pgfplots.fillbetween}\usetikzlibrary{intersections, pgfplots.fillbetween}
\rowcolors{2}{miugray!30}{miugray!10}

\usetheme{Madrid}  %% Themenwahl
\useoutertheme{split} % Alternatively: miniframes, infolines, split
%\useinnertheme{circles}


\setbeamertemplate{itemize item}{\color{miugray}$\blacksquare$}
\setbeamertemplate{itemize subitem}{\color{uniblau}$\bullet$}



%\newcommand{\metermilli}{\mskip3mu \text{mm}}
\newcommand{\cm}{\mskip3mu \text{cm}}
\newcommand{\m}{ \text{m}}
\newcommand{\lite}{ \text{l}}
\newcommand{\kg}{ \text{kg}}
\newcommand{\g}{ \text{g}}
\newcommand{\km}{ \text{km}}
\newcommand{\h}{ \text{h}}
\newcommand{\s}{ \text{s}}
\newcommand{\tonne}{ \text{t}}


\newcommand{\RA}{$\rightarrow$~}
\newcommand{\pp}{\par \vspace*{2mm}}
\newcommand{\ppbig}{\par \vspace*{4mm}}

\definecolor{miudiutex}{HTML}{1d2926}
\definecolor{miugray}{HTML}{b4cccf}
\definecolor{miugray2}{HTML}{4d7365}
\definecolor{white2}{HTML}{f5f7f8}
\definecolor{red}{HTML}{cf1f5c}
\definecolor{green}{HTML}{00A36C}
\definecolor{delphinidin}{HTML}{315794}
\definecolor{uniblau}{HTML}{004e9f}
\definecolor{unigelb}{HTML}{fcba00}
\definecolor{unigrau}{HTML}{909085}
\definecolor{pelargonidin}{HTML}{b33807}
\definecolor{cyanidin}{HTML}{ba0746}
\definecolor{paeonidin}{HTML}{d15ca6}
\definecolor{petunidin}{HTML}{b56fed}
\definecolor{malvidin}{HTML}{8a2d8c}

%\setbeamercolor{palette primary}{bg=pelargonidin!35,fg=white}
%\setbeamercolor{palette secondary}{bg=delphinidin!45,fg=white}
\setbeamercolor{palette primary}{bg=uniblau!70,fg=white}
\setbeamercolor{palette secondary}{bg=uniblau!45,fg=white}
\setbeamercolor{palette tertiary}{bg=unigelb!70,fg=unigrau}
\setbeamercolor{palette quaternary}{bg=miugray,fg=white}
\setbeamercolor{structure}{fg=black} % itemize, enumerate, etc
\setbeamercolor{section in toc}{fg=red} % TOC sections

\setbeamersize{text margin left=1cm,text margin right=1cm}


\useinnertheme{tcolorbox} 
\newcommand{\goodtk}[2]{
\begin{tcolorbox}[
	enhanced,
	colframe=uniblau!75,
	sharp corners,
	boxsep=5pt,
	title=\bfseries \sffamily #1,
	colback=white,
	coltitle=black,
	attach boxed title to top text left={yshift=-\tcboxedtitleheight/2},
	boxed title style={colback=white,colframe=white}
]
	#2
	\end{tcolorbox}
}
	
\newcommand{\antwort}[1]{
\begin{tcolorbox}[
	enhanced,
	colframe=unigrau!65,
	sharp corners,
	boxsep=1pt,
	colback=white,
	coltitle=black,
]
	#1
	\end{tcolorbox}
}


\title{7. Übung}
\author{PD Dr. Helene Loos}
\date{\today}

\begin{document}
%\maketitle
\small



\begin{frame}
\begin{center}

\LARGE Übung 7  \par \normalsize
Prozessbezogene Lebensmitteltechnologie
\end{center}
\small
\end{frame}



\begin{frame}
\begin{minipage}{.4\textwidth}


{\color{miugray}$\blacksquare ~$} \textbf{Aufgabe 1} \par \vspace*{2mm}

Wasser siedet unter 1 atm (101,325 kPa) Druck bei 100 °C. Bei welcher Temperatur siedet
eine 0,1 molare wässrige Lösung aus Kochsalz? Die Dichte von Wasser soll als 1.000 kg/m$^3$ angenommen werden, die spezifische Verdampfungsenthalpie beträgt 2.260 kJ/kg. 


\end{minipage}%
\hspace*{4mm}
\begin{minipage}{.6\textwidth}
\antwort{
\footnotesize
\uncover<2->{
Formel zur Bearbeitung der Aufgabe:}
\uncover<3->{\begin{align*}
\Delta T \approx \frac{R \cdot T_L^2}{\rho_L \cdot \Delta h_v} \cdot c_i
\end{align*}}
\uncover<4->{Werte einsetzen:}
\uncover<5->{\begin{align*}
\Delta T &\approx \frac{8,314 ~\text{J/K mol} \cdot (373,15 ~\text{K})^2}{1.000 ~\text{kg/m$^3$} \cdot 2,26 \cdot 10^6 ~\text{J/kg}} \cdot 0,2 \cdot 10^{3}~\text{mol/m}^3 \\
&\approx 0,1 ~\text{K}
\end{align*}}
}
\end{minipage}%


\end{frame}



\begin{frame}
\begin{minipage}{.38\textwidth}


{\color{miugray}$\blacksquare ~$} \textbf{Aufgabe 2} \par \vspace*{2mm}
Ein Verdampfer ist 1,6 m hoch und enthält siedendes Wasser bei 100°C. Um wieviel erhöht
sich die Siedetemperatur unten im Verdampfer, weil dort wegen des Gewichts der
Wassersäule ein höherer Druck herrscht? Vernachlässigen Sie den Dampfanteil der
siedenden Flüssigkeit. Die Dichte von Wasser bei 100 °C soll als 960 kg/m$^3$
 angenommen werden, die spezifische Verdampfungsenthalpie beträgt 2.260 kJ/kg und die molare Masse
von Wasser 18 g/mol. Als Umgebungsdruck sollen 100 kPa angenommen werden. 




\end{minipage}%
\hspace*{4mm}
\begin{minipage}{.62\textwidth}
\antwort{
\footnotesize
\uncover<2->{Formel zur Bearbeitung der Aufgabe:}
\uncover<3->{\begin{align*}
\Delta T = \frac{\rho_{lg} \cdot g \cdot h \cdot R \cdot T_L^2}{\Delta h_{v} \cdot M_L \cdot p}
\end{align*}}
\uncover<4->{Werte einsetzen:}
\uncover<5->{\begin{align*}
\Delta T &= \frac{960 ~\text{kg/m}^3 \cdot 9,81 ~\text{m/s}^2 \cdot 1,6 ~\text{m} \cdot 8,314 ~\text{J/mol K} \cdot (373,15 ~\text{K})^2}{2,26 \cdot 10^6 ~\text{J/kg} \cdot 0,018 ~\text{kg/mol} \cdot 100.000 ~\text{Pa}} \\
&= 4,29 ~\text{K} 
\end{align*}}
}
\end{minipage}%


\end{frame}














\begin{frame}
\begin{minipage}{.3\textwidth}


{\color{miugray}$\blacksquare ~$} \textbf{Aufgabe 3} \par \vspace*{2mm}
Ein Lösungsmittel hat einen Dampfdruck von 0,126 bar bei 25 °C und einen Dampfdruck von
1,013 bar bei 77°C. Wie groß ist der Dampfdruck bei 51°C?



\end{minipage}%
\hspace*{4mm}
\begin{minipage}{.7\textwidth}
\antwort{
\footnotesize
\uncover<2->{
Wir wissen, dass
\begin{align*}
\ln(\frac{p_2}{p_1}) = \ln(\frac{p_3}{p_1})
\end{align*}}
\uncover<3->{Clausius-Clapeyron:
\begin{align*}
\ln(\frac{p_2}{p_1}) = \frac{\Delta h_{vM}}{R} \cdot \left( \frac{1}{T_1} - \frac{1}{T_2} \right)
\end{align*}}
\uncover<4->{Umformen:}
\uncover<5->{
\begin{align*}
\Delta h_{vM} = \frac{R \cdot T_1 \cdot T_2 \cdot \ln(\frac{p_2}{p_1})}{T_2 - T_1}
\end{align*}}
\uncover<6->{Werte einsetzen:}
\uncover<7->{
\begin{align*}
\Delta h_{vM} &= \frac{8,314 ~\text{J/(mol K)} \cdot 298,15 ~\text{K} \cdot 350,15 ~\text{K} \cdot \ln(\frac{1,013 \cdot 10^5 ~\text{Pa}}{0,126 \cdot 10^5 ~\text{Pa}})}{350,15 ~\text{K} - 298,15 ~\text{K}} \\
 &= 34.791,6 ~\text{J/mol}
\end{align*}}

}
\end{minipage}%


\end{frame}
\begin{frame}
\begin{minipage}{.3\textwidth}


{\color{miugray}$\blacksquare ~$} \textbf{Aufgabe 3} \par \vspace*{2mm}
Ein Lösungsmittel hat einen Dampfdruck von 0,126 bar bei 25 °C und einen Dampfdruck von
1,013 bar bei 77°C. Wie groß ist der Dampfdruck bei 51°C?



\end{minipage}%
\hspace*{4mm}
\begin{minipage}{.7\textwidth}
\antwort{
\footnotesize
\uncover<2->{
Clausius-Clapeyron nach $p_3$ umstellen:
}
\begin{align*}
\uncover<3->{\ln(p_3) &= \frac{\Delta h_{vM}}{R} \cdot \left( \frac{1}{T_1} - \frac{1}{T_3} \right) + \ln(p_1) \\}
\uncover<4->{\ln(p_3) &= \frac{34.791,6 ~\text{J/mol}}{8,314 ~\text{J/(mol K)}} \cdot \left( \frac{1}{298,15 ~\text{K}} - \frac{1}{324,15 ~\text{K}} \right) + \ln(0,126 \cdot 10^5 ~\text{Pa}) \\}
\uncover<5->{&= 10,57 ~\text{Pa}\qquad \qquad| e^x \\ \\}
\uncover<6->{p_3 &= 38.948,67 ~\text{Pa} \\
&= 0,389 ~\text{bar}}
\end{align*}

}
\end{minipage}%


\end{frame}























\begin{frame}
\begin{minipage}{.28\textwidth}


{\color{miugray}$\blacksquare ~$} \textbf{Aufgabe 4} \par \vspace*{2mm}
Ihnen liegt ein Behälter mit einer Flüssigkeit vor. Diese Flüssigkeit besitzt bei einer
Temperatur von 25 °C einen Dampfdruck von 1 atm (101,325 kPa). Berechnen Sie den
Dampfdruck, wenn die Temperatur in dem Behälter um 80 °C steigt. Die molspezifische
Verdampfungsenthalpie der Flüssigkeit beträgt 40,657 kJ/mol am Siedepunkt. 


\end{minipage}%
\hspace*{4mm}
\begin{minipage}{.72\textwidth}
\antwort{
\footnotesize
\begin{align*}
\uncover<2->{\ln(p_3) &= \frac{40.657 ~\text{J/mol}}{8,314 ~\text{J / (mol K)}} \cdot \left( \frac{1}{298,15~\text{K}} - \frac{1}{378,15~\text{K}} \right) + \ln(101,325 \cdot 10^3 ~\text{Pa}) \\}
\uncover<3->{&= 14,99~\text{Pa}}\uncover<4->{ \qquad \qquad| e^x \\ \\}
\uncover<5->{p_3 &= 3.255.914,7 ~\text{Pa} \\}
\uncover<6->{&= 32,5 ~\text{bar}}
\end{align*}

}

\end{minipage}%



\end{frame}















\begin{frame}
\begin{minipage}{.28\textwidth}


{\color{miugray}$\blacksquare ~$} \textbf{Aufgabe 5} \par \vspace*{2mm}
Wasser siedet unter 1 atm (101,325 kPa) Druck bei 100 °C. Bei welcher Temperatur siedet
Wasser in einem Dampfkochtopf, welcher 2 bar Druck hat? Die spezifische
Verdampfungsenthalpie beträgt 2260 kJ/kg, das Molgewicht 18 g/mol.



\end{minipage}%
\hspace*{4mm}
\begin{minipage}{.72\textwidth}

\antwort{
\footnotesize
\begin{align*}
\uncover<2->{\ln(\frac{p_2}{p_1}) \cdot \frac{R}{\Delta h_{vM}}+\frac{1}{T_2}&=\frac{1}{T_1} \\}
\uncover<3->{\ln(\frac{101.325~\text{Pa}}{200.000~\text{Pa}}) \cdot \frac{8,314 ~\text{J/(mol K)}}{2,26 \cdot 10^6 ~\text{J/kg} \cdot 0,018 ~\text{kg/mol}}+\frac{1}{373,15 ~\text{K}}&=
\frac{2,541 \cdot 10^{-3}}{\text{K}} \\}
\uncover<4->{&=393,55 ~\text{K}}
\end{align*} 

}
\end{minipage}%
\end{frame}


















\begin{frame}


\begin{minipage}{.28\textwidth}


{\color{miugray}$\blacksquare ~$} \textbf{Aufgabe 6 a)} \par \vspace*{2mm}
Wasser siedet unter 1 atm (101,325 kPa) Druck bei 100 °C.
\pp
Bei welcher Temperatur siedet Wasser am Mount Everest (0,325 bar)? Die spezifische
Verdampfungsenthalpie beträgt 2.260 kJ/kg, das Molgewicht 18 g/mol.
 


\end{minipage}%
\hspace*{4mm}
\begin{minipage}{.72\textwidth}

\antwort{
\footnotesize
\begin{align*}
\uncover<2->{\ln(\frac{p_2}{p_1})\cdot \frac{R}{\Delta h_{vM}}+\frac{1}{T_2}&=\frac{1}{T_1} \\}
\uncover<3->{\ln(\frac{101.325~\text{Pa}}{32.500~\text{Pa}}) \cdot \frac{8,314 ~\text{J/(mol K)}}{2,26 \cdot 10^6 ~\text{J/kg} \cdot 0,018 ~\text{kg/mol}}+\frac{1}{373,15~\text{K}}&=\frac{2,915 \cdot 10^{-3}}{\text{K}}\\ }
\uncover<4->{&=343 ~\text{K}}
\end{align*} 
}





\end{minipage}%
\end{frame}


\begin{frame}
\begin{minipage}{.5\textwidth}


{\color{miugray}$\blacksquare ~$} \textbf{Aufgabe 6 b)} \par \vspace*{2mm}
Für einen Erhitzungsprozess wählen Sie bei 100°C etwa 8 Minuten. Wie lange dauert der Erhitzungsprozess am Mount Everest, wenn der z-Wert 20 K beträgt? Nehmen Sie an, dass
das Ziel des Erhitzungsprozesses auch bei der reduzierten Temperatur erreicht wird. 

\end{minipage}%
\hspace*{4mm}
\begin{minipage}{.5\textwidth}
\antwort{

\footnotesize
\begin{align*}
\uncover<2->{F_1 &= F_2 \cdot 10^{\frac{T_2-T_1}{z}} \\}
\uncover<3->{F_1 &= 8~\text{min} \cdot 10^{\frac{373,15 ~\text{K}-343 ~\text{K}}{20 ~\text{K}}} \\}
\uncover<4->{ &= 257,38 ~\text{min}}
\end{align*}

}
\end{minipage}%
\end{frame}



\end{document}
